\documentclass[11pt]{article}
\usepackage[margin=0.8in]{geometry}
\usepackage[T1]{fontenc}
\usepackage{lmodern}
\usepackage{microtype}
\usepackage{xcolor}
\usepackage{enumitem}
\usepackage{booktabs}
\usepackage{hyperref}
\hypersetup{colorlinks=true,linkcolor=blue,urlcolor=blue,pdftitle={Week 7 Guided Path: Contract and Swap}}
\setlist{nosep,leftmargin=*}
\definecolor{navy}{HTML}{17365D}
\definecolor{softblue}{HTML}{EAF2FB}
\newcommand{\code}[1]{\texttt{#1}}
\newcommand{\heading}[1]{\vspace{0.6em}\noindent\colorbox{softblue}{\parbox{\dimexpr\linewidth-2\fboxsep}{\large\bfseries\color{navy}#1}}\vspace{0.25em}}
\begin{document}
\begin{center}
{\LARGE\bfseries Week 7 Guided Path: Contract and Swap}\\[0.35em]
{\large COMSC-1053 \quad|\quad Reasoning Odyssey Gate}
\end{center}

\noindent\textbf{The goal.} Your program already depends on other objects.
This week you will find one real dependency, write down its promise as an
explicit abstract base class (ABC), plug in two different collaborators, and
prove the caller can work with either one. You are not building a second
project. You are making a small, honest improvement to the world you built in
Weeks 4--6.

\heading{What you will make}
\begin{enumerate}
  \item An \textbf{ABC} with one abstract operation that states the collaborator's promise.
  \item Two concrete collaborators that inherit from the ABC, preserve meaningful state, and behave differently.
  \item One caller that depends on the promise instead of branching on concrete types.
  \item One focused \texttt{unittest} that uses the same caller and request with both collaborators.
  \item A short World Bible update explaining the change, evidence, and remaining debt.
\end{enumerate}

\heading{First, see the complete example}
The numbered Frontier Settlement lesson is one guided example. It is not a
template to copy. Its \code{ExpeditionPlanner} asks a \code{SupplySource} for a
\code{SupplyQuote}. A \code{SettlementWarehouse} and a \code{TradingPost} both
provide quotes, but their rules differ: the trading post protects a reserve.

\begin{center}
\begin{tabular}{@{}r l p{3.45in}@{}}
\toprule
\textbf{Step} & \textbf{File} & \textbf{Question it answers}\\
\midrule
1 & \code{step\_01\_values.py} & What values and errors cross this boundary?\\
2 & \code{step\_02\_contracts.py} & What must every supply source promise?\\
3 & \code{step\_03\_validation.py} & What inputs are valid before delegation?\\
4 & \code{step\_04\_warehouse.py} & How can one source keep meaningful state?\\
5 & \code{step\_05\_trading\_post.py} & How can another source follow a different rule?\\
6 & \code{step\_06\_planner.py} & What work belongs to the caller boundary?\\
7 & \code{step\_07\_tests.py} & What evidence proves the shared promise?\\
8 & \code{step\_08\_demo.py} & What does a visible swap look like?\\
9 & \code{step\_09\_world\_bible.py} & What changed and what debt remains?\\
\bottomrule
\end{tabular}
\end{center}

\noindent Run it from the Week 7 directory:
\begin{verbatim}
python class_example.py test
python class_example.py demo
python class_example.py bible
\end{verbatim}
The first command runs 34 tests. The test file is deliberately named
\code{step\_07\_tests.py}, so use this entry point rather than generic test
discovery at the directory root.

\heading{The idea in plain language}
An explicit contract tells the caller what it may trust. In the example, the
caller may trust that a supply source can answer \code{quote\_supply(...)} with
a valid \code{SupplyQuote}. It may \emph{not} assume the source is a warehouse,
a trading post, or any future kind of source.

The caller boundary is the moment \code{ExpeditionPlanner} hands a request to
the collaborator and accepts a result back. A good boundary does three jobs:
\begin{itemize}
  \item validates request data before it crosses out;
  \item checks the returned value still satisfies the promise; and
  \item adds useful context if the collaborator fails.
\end{itemize}
This is why the caller has no warehouse-specific or trading-post-specific
branch. Swapping a collaborator changes behavior without rewriting the caller.

\heading{Build your own version}
\begin{enumerate}
  \item Look back at your Weeks 4--6 world. Find a dependency that already
  exists in the program. Ask: What does my caller need from this object?
  \item State that need as a narrow promise. Prefer one operation with a
  meaningful return value over a broad interface.
  \item Create an \code{ABC} and mark the operation with \code{@abstractmethod}.
  \item Create two concrete subclasses. Give each its own state and rule so a
  substitution changes behavior for a real reason.
  \item Update the caller to receive the collaborator and use only the contract.
  \item Write the focused swap test. Use the same caller operation and input for
  each collaborator; assert only what the contract promises.
  \item Run the tests, demonstrate the result, and record the design decision in
  your World Bible.
\end{enumerate}

\heading{A focused test looks for a promise, not a sentence}
A brittle test checks a warehouse's exact wording. A contract test checks the
result that every valid collaborator promises. The key shape is:
\begin{verbatim}
for source in [warehouse, trading_post]:
    planner = ExpeditionPlanner(source)
    decision = planner.evaluate_supply_request(...)
    self.assertIsInstance(decision, SupplyDecision)
\end{verbatim}
Add collaborator-specific tests separately when you need to prove a particular
rule, such as the trading post's reserve. That keeps the shared contract test
honest and keeps your failures easy to understand.

\heading{Your graded discussion post (40 points)}
Your post is a shared learning resource: classmates can see it. Link your
repository or attach the code, then write a concise explanation with:
\begin{center}
\begin{tabular}{@{}p{1.6in}r p{3.45in}@{}}
\toprule
\textbf{Evidence} & \textbf{Points} & \textbf{What a reader should find}\\
\midrule
Real contract and two collaborators & 15 & An ABC, abstract operation, and two stateful implementations with genuinely different behavior.\\
Focused swap test & 10 & The same caller and request work with both collaborators; the assertion matches the shared promise.\\
Caller-boundary explanation & 10 & Plain-language explanation of validation, delegation, returned-value checking, and failure context.\\
Demonstration and World Bible & 5 & A runnable result plus what changed, evidence used, and remaining debt.\\
\bottomrule
\end{tabular}
\end{center}
Then reply constructively to at least two classmates. Refer to a specific
design or test, ask a useful question, or name an idea you can carry into your
own program. Do not post private information about yourself or anyone else.

\heading{Optional reference worlds}
The folders beginning \code{contract\_swap\_example} and the
\code{week\_07\_contract\_and\_swap} folder are optional, self-contained
illustrations in other settings. Read one only if it helps. They are neither
extra assignments nor code to copy. The numbered Frontier Settlement path is
the one complete walkthrough.

\vfill
\noindent\textit{A strong Week 7 submission is small, specific, tested, and
honest about what it still does not solve.}
\end{document}
